\Section{Logic Gates}
\begin{multicols}{2}
\textbf{NOT Gate (Inverter)}\\
NOT gate accepts a single input and outputs the opposite value.
NOT gate is also known as an “INVERTER” as it inverts the input value.


\begin{center}
\vspace{2ex}
\begin{circuitikz} \draw
   (0,2) node[not port] (not) {}
   (not.in) -- ++(-0.25,0) node[left] {$A$}
   (not.out) -- ++(0.25,0) node[right] {$R$};
\end{circuitikz}

\vspace{2ex}
\begin{tabularx}{5cm}{|C|C|} \hline
   A & R \\\hline
   0 & 1 \\\hline
   1 & 0 \\\hline
\end{tabularx}

% \vspace{2ex}
% \begin{circuitikz}[x=0.5cm,y=0.5cm]
%    \draw
%    (0,0) |- 
%    (2,3) |- (3,3.75) to[nos,o-o] (4,3.75) -| (5,3)
%    (2,3) |- (3,2.25) to[nos,o-o] (4,2.25) -| (5,3)
%    (5,3) -- (7,3) to[led] (7,0) to[battery] (0,0)
%    ;
% \end{circuitikz}
\end{center}


\vspace{4ex}
\textbf{AND Gate}\\
AND Operation produces HIGH output (1) only if all the inputs to the digital circuit are HIGH (1). Otherwise output will be LOW (0).

\begin{center}
\vspace{2ex}
\begin{tabularx}{5cm}{|C|C|C|} \hline
   A & B & R \\\hline
   0 & 0 & 0 \\\hline
   0 & 1 & 0 \\\hline
   1 & 0 & 0 \\\hline
   1 & 1 & 1 \\\hline
\end{tabularx}

\begin{circuitikz} \draw
   (0,2) node[and port] (and) {}
   (and.in 1) -- ++(-0.25,0) node[left] {$A$}
   (and.in 2) -- ++(-0.25,0) node[left] {$B$}
   (and.out) -- ++(0.25,0) node[right] {$R$};
\end{circuitikz}
\end{center}


\vspace{4ex}
\textbf{OR Gate}\\
OR Operation produces HIGH output (1) if at least one of the inputs is HIGH (1).

\begin{center}
\vspace{2ex}
\begin{tabularx}{5cm}{|C|C|C|} \hline
   A & B & R \\\hline
   0 & 0 & 0 \\\hline
   0 & 1 & 1 \\\hline
   1 & 0 & 1 \\\hline
   1 & 1 & 1 \\\hline
\end{tabularx}

\begin{circuitikz} \draw
   (0,2) node[or port] (or) {}
   (or.in 1) -- ++(-0.25,0) node[left] {$A$}
   (or.in 2) -- ++(-0.25,0) node[left] {$B$}
   (or.out) -- ++(0.25,0) node[right] {$R$};
\end{circuitikz}
\end{center}


\vspace{4ex}
\textbf{NAND Gate}\\
NAND Operation produces LOW output (0) only if all the inputs to the digital circuit are HIGH (1). Otherwise output will be HIGH (1).

\begin{center}
\vspace{2ex}
\begin{tabularx}{5cm}{|C|C|C|} \hline
   A & B & R \\\hline
   0 & 0 & 1 \\\hline
   0 & 1 & 1 \\\hline
   1 & 0 & 1 \\\hline
   1 & 1 & 0 \\\hline
\end{tabularx}

\begin{circuitikz} \draw
   (0,2) node[nand port] (nand) {}
   (nand.in 1) -- ++(-0.25,0) node[left] {$A$}
   (nand.in 2) -- ++(-0.25,0) node[left] {$B$}
   (nand.out) -- ++(0.25,0) node[right] {$R$};
\end{circuitikz}
\end{center}


\vspace{4ex}
\textbf{NOR Gate}\\
NOR Operation produces LOW output (0) if at least one of the inputs is HIGH (1).


\begin{center}
\vspace{2ex}
\begin{tabularx}{5cm}{|C|C|C|} \hline
   A & B & R \\\hline
   0 & 0 & 1 \\\hline
   0 & 1 & 0 \\\hline
   1 & 0 & 0 \\\hline
   1 & 1 & 0 \\\hline
\end{tabularx}

\begin{circuitikz} \draw
   (0,2) node[nor port] (nor) {}
   (nor.in 1) -- ++(-0.25,0) node[left] {$A$}
   (nor.in 2) -- ++(-0.25,0) node[left] {$B$}
   (nor.out) -- ++(0.25,0) node[right] {$R$};
\end{circuitikz}
\end{center}


\vspace{4ex}
\textbf{XOR Gate}\\
XOR Operation produces LOW output (0) if both of the inputs are same and produces HIGH output (1) if the inputs are different.

\begin{center}
\vspace{2ex}
\begin{tabularx}{5cm}{|C|C|C|} \hline
   A & B & R \\\hline
   0 & 0 & 0 \\\hline
   0 & 1 & 1 \\\hline
   1 & 0 & 1 \\\hline
   1 & 1 & 0 \\\hline
\end{tabularx}

\begin{circuitikz} \draw
   (0,2) node[xor port] (xor) {}
   (xor.in 1) -- ++(-0.25,0) node[left] {$A$}
   (xor.in 2) -- ++(-0.25,0) node[left] {$B$}
   (xor.out) -- ++(0.25,0) node[right] {$R$};
\end{circuitikz}
\end{center}


\vspace{4ex}
\textbf{XNOR Gate}\\
XNOR Operation produces HIGH output (1) if both of the inputs are same and produces LOW output (0) if the inputs are different.

\begin{center}
\vspace{2ex}
\begin{tabularx}{5cm}{|C|C|C|} \hline
   A & B & R \\\hline
   0 & 0 & 1 \\\hline
   0 & 1 & 0 \\\hline
   1 & 0 & 0 \\\hline
   1 & 1 & 1 \\\hline
\end{tabularx}

\begin{circuitikz}
   \draw
   (0,2) node[xnor port] (xnor) {}
   (xnor.in 1) -- ++(-0.25,0) node[left] {$A$}
   (xnor.in 2) -- ++(-0.25,0) node[left] {$B$}
   (xnor.out) -- ++(0.25,0) node[right] {$R$};
\end{circuitikz}
\end{center}
\end{multicols}